\chapter*{Introduction Générale}
\addcontentsline{toc}{chapter}{Introduction Générale}

Dans un contexte de transformation numérique accélérée, la gestion des ressources
humaines et de la paie constitue un levier stratégique pour toute organisation. En
Tunisie, ce domaine est régi par un cadre juridique dense et évolutif~: cotisations
CNSS, barème progressif de l'IRPP, taux CAVIS, Contribution Sociale de Solidarité et
dispositions annuelles de la Loi de Finances~\cite{loifinance}. Garantir simultanément
la conformité légale et la maîtrise de la charge administrative représente un défi
structurel permanent pour les PME et les cabinets d'expertise comptable tunisiens.

L'analyse du marché révèle une coexistence préoccupante entre des pratiques manuelles
fondées sur Excel et des logiciels de bureau isolés aux fonctionnalités figées, dépourvus
de mobilité, d'architecture multi-tenant et d'intelligence artificielle. Ces lacunes
engendrent des risques de non-conformité et une inefficacité opérationnelle qui pèsent
directement sur la compétitivité des entreprises.

C'est dans ce cadre que s'inscrit notre projet de fin d'études, réalisé au sein de
\textbf{Grow Up Technologies}, visant à concevoir et développer \textbf{PaieZone~RH}~---
une plateforme \textit{SaaS} multi-tenant de gestion RH et de paie intégrée à
\textit{ComptaZone}. Le projet poursuit quatre objectifs~: automatiser le cycle RH
complet, garantir la conformité à la Loi de Finances 2026 via un moteur de calcul
paramétrable, assurer l'isolation des données par schéma PostgreSQL dédié par tenant,
et déployer \textbf{PaieBot}~--- un assistant conversationnel RH local basé sur
Ollama/phi3, combinant un pipeline Text-to-SQL sécurisé et une architecture
RAG~\cite{rag} pour les questions réglementaires.

Le développement a suivi la méthodologie \textbf{Scrum}~\cite{scrum}, structuré en
quatre sprints de trois semaines couvrant 61 User Stories réparties en seize épics.
La conception s'appuie sur \textbf{UML}~\cite{uml} et l'architecture \textit{N-tiers}
de \textbf{JHipster~9.x}~\cite{jhipster} (\textbf{Spring Boot~4.0}~\cite{springboot}
+ \textbf{Angular~18}~\cite{angular}). Ce rapport s'organise en six chapitres~:
le premier expose le contexte général~; le second couvre la planification et
l'architecture~; les chapitres trois à six détaillent successivement la réalisation
de chaque sprint, depuis le socle technique jusqu'au déploiement de l'assistant IA.
